\documentclass[11pt]{article}

\usepackage[margin=1in]{geometry}
\usepackage{amsmath}
\usepackage{amssymb}
\usepackage{amsthm}

\newtheorem{theorem}{Theorem}
\newtheorem{lemma}[theorem]{Lemma}
\newtheorem{proposition}[theorem]{Proposition}
\newtheorem{corollary}[theorem]{Corollary}
\theoremstyle{definition}
\newtheorem{conjecture}[theorem]{Conjecture}
\theoremstyle{remark}
\newtheorem{remark}[theorem]{Remark}

\newcommand{\R}{\mathbb{R}}
\newcommand{\Z}{\mathbb{Z}}
\newcommand{\C}{\mathbb{C}}
\newcommand{\T}{\mathbb{T}}
\newcommand{\Lam}{\Lambda}
\newcommand{\rtwo}{r_2}

\title{Refutation of Conjecture 00000000277:\\
the maximum Laplace multiplicity on a flat torus is not $6$}
\author{earthking11}
\date{13 September 2026}

\begin{document}
\maketitle

\begin{abstract}
We disprove Conjecture 00000000277.  The conjecture asserts that $M$, the
maximum multiplicity of a Laplace eigenvalue on a planar flat torus (supremum
over all such tori), equals $6$, and that multiplicity six is attained only on
a lattice torus whose automorphism group has order thirty-two.  Under the
literal reading of $M$ (maximum over \emph{all} eigenvalues) the claim fails
outright: on the square torus $\R^2/\Z^2$ the eigenvalue $4\pi^2 N$ has
multiplicity $\rtwo(N)$, the number of representations of $N$ as a sum of two
squares, and $\rtwo(5^k)=4(k+1)\to\infty$, so the multiplicity is unbounded and
$M=\infty$.  Independently, the order-thirty-two clause is impossible under
\emph{either} reading: the crystallographic restriction forces every
lattice-preserving rotation to have order in $\{1,2,3,4,6\}$, so the
automorphism group of a planar lattice has order at most $12$.  We are careful
to record that if $M$ had been intended as the multiplicity of the
\emph{first non-zero} eigenvalue only, then the value $6$ is correct (attained
by the hexagonal lattice); nevertheless the statement as written is false.
\end{abstract}

\section{The conjecture}

\begin{conjecture}[00000000277]
$M(T)$ is the maximum multiplicity of a Laplace eigenvalue on a planar flat
torus $T=\C/\Lam$, and $M$ is its supremum over all flat tori.  Then
$M=6$ (six is known to be attainable, and upper bounds $\ge 6$ are constrained
by congruence structure); moreover multiplicity six is attained only on the
lattice torus with maximal automorphism group of order thirty-two.
\end{conjecture}

We show that this statement is false.  Section~\ref{sec:spectral} sets up the
spectrum of a flat torus, Section~\ref{sec:r2} produces unbounded
multiplicities on the square torus, Section~\ref{sec:cryst} rules out the
order-thirty-two clause, and Section~\ref{sec:honest} states precisely the
sense in which the number $6$ is nevertheless correct.

\section{The spectrum of a flat torus}
\label{sec:spectral}

Let $\Lam\subset\R^2\cong\C$ be a lattice of full rank and let
$T=\R^2/\Lam$ be the associated flat torus, equipped with the flat metric
induced by the Euclidean metric on $\R^2$.  The Laplace--Beltrami operator
$-\Delta_T$ has spectrum
\begin{equation}
\label{eq:spectrum}
\{4\pi^2|k|^2 \;:\; k\in\Lam^{*}\},
\qquad
\Lam^{*}=\{\,k\in\R^2 : \langle k,\lambda\rangle\in\Z
\text{ for all }\lambda\in\Lam\,\},
\end{equation}
where $\Lam^{*}$ is the dual lattice, and the multiplicity of the eigenvalue
$4\pi^2|k|^2$ is the number of dual-lattice vectors of that squared length:
\begin{equation}
\label{eq:mult}
\mathrm{mult}_T(4\pi^2 N)=\#\{\,k\in\Lam^{*} : |k|^2=N\,\}.
\end{equation}
This is the standard Fourier expansion: the functions
$x\mapsto e^{2\pi i\langle k,x\rangle}$ for $k\in\Lam^{*}$ form a complete
orthogonal system of eigenfunctions with eigenvalues $4\pi^2|k|^2$.  In
particular each multiplicity is finite and even whenever the length is
non-zero, because $k$ and $-k$ are distinct dual vectors of equal length.

For the square torus $\Lam=\Z^2=\Z[i]$, the lattice is self-dual,
$\Lam^{*}=\Z^2$, and the squared lengths are exactly the integers representable
as a sum of two squares:
\begin{equation}
\label{eq:square-mult}
\mathrm{mult}_{\Z[i]}(4\pi^2 N)
=\#\{\,(a,b)\in\Z^2 : a^2+b^2=N\,\}
=:\rtwo(N).
\end{equation}

\section{Unbounded multiplicity on the square torus}
\label{sec:r2}

The arithmetic function $\rtwo(N)$ is classical.  Since $\Z[i]$ is a unique
factorisation domain, the number of representations of $N$ as a sum of two
squares is
\begin{equation}
\label{eq:r2-formula}
\rtwo(N)=4\,(d_1(N)-d_3(N)),
\end{equation}
where $d_1(N)$ (resp.\ $d_3(N)$) counts the positive divisors of $N$ congruent
to $1$ (resp.\ $3$) modulo $4$; this is Jacobi's two-square theorem.  We will
only need the case $N=5^k$, whose divisors are $1,5,5^2,\dots,5^k$, all
congruent to $1 \bmod 4$.  Hence
\begin{equation}
\label{eq:r2-5k}
\rtwo(5^k)=4(k+1) \xrightarrow[k\to\infty]{} \infty .
\end{equation}

The first values are
\begin{equation}
\rtwo(1)=4,\quad
\rtwo(5)=8,\quad
\rtwo(25)=12,\quad
\rtwo(125)=16,\quad
\rtwo(625)=20,\quad
\rtwo(3125)=24 .
\end{equation}
Concretely $\rtwo(25)=12$ comes from
$(a,b)\in\{(\pm5,0),(0,\pm5),(\pm3,\pm4),(\pm4,\pm3)\}$, and in general
$\rtwo(5^k)=4(k+1)$ counts the pairs $(\pm 2^i 5^{j},\pm 2^{i'}5^{j'})$ with
$i+i'=j+j'=k$.

\begin{proposition}[Unbounded multiplicity; $M=\infty$]
\label{prop:unbounded}
For every $m$ there is a planar flat torus carrying a Laplace eigenvalue of
multiplicity at least $m$.  Consequently the supremum $M$ of all such
multiplicities is $+\infty$, and in particular $M\neq6$.
\end{proposition}

\begin{proof}
Take the square torus $T=\R^2/\Z^2$ and $N=5^k$ for $k\ge \lceil m/4\rceil-1$.
By \eqref{eq:square-mult} the eigenvalue $4\pi^2 N$ has multiplicity
$\rtwo(5^k)=4(k+1)\ge m$.  Thus multiplicities are unbounded, so the supremum
is infinite.  Already for $k=2$, i.e.\ $N=25$, the multiplicity is
$\rtwo(25)=12>6$, so the claimed value $M=6$ fails on a single eigenvalue.
\end{proof}

\begin{remark}
Proposition~\ref{prop:unbounded} addresses the literal reading of the
definition, in which $M(T)$ is the maximum multiplicity over all eigenvalues
of $-\Delta_T$.  The conclusion $M=\infty$ is not a consequence of some exotic
lattice: the square torus, the simplest and most classical flat torus, already
witnesses multiplicities $12,16,20,24,\dots$.
\end{remark}

\section{The order-thirty-two lattice is impossible}
\label{sec:cryst}

The second clause of the conjecture asserts that multiplicity six is attained
only on a lattice torus whose automorphism group has order thirty-two.  This
clause is false under either reading of $M$, for a reason independent of the
spectral computation above.

Let $\Lam\subset\R^2$ be a lattice and let
$\mathrm{Aut}(\Lam)=\{g\in O(2) : g\Lam=\Lam\}$ be its orthogonal automorphism
group, a finite group.  Every $g\in\mathrm{Aut}(\Lam)$ is a rotation or a
reflection; we consider the rotations.

\begin{theorem}[Crystallographic restriction]
\label{thm:cryst}
If $g\in O(2)$ has finite order $n$ and preserves a full-rank lattice
$\Lam\subset\R^2$, then the trace $\operatorname{tr}(g)=2\cos(2\pi/n)$ is an
integer, and therefore
\begin{equation}
n\in\{1,2,3,4,6\}.
\end{equation}
\end{theorem}

\begin{proof}
In an integral basis of $\Lam$, the map $g$ is represented by a matrix
$A\in GL_2(\Z)$.  Since $g$ has order $n$, so does $A$, hence $A^n=I$.  The
trace $\operatorname{tr}(A)$ is an integer, and it equals the sum of the two
eigenvalues of $A$ over $\C$.  Those eigenvalues are roots of unity
$\zeta,\zeta^{-1}$ with $\zeta^n=1$, so for a rotation of order $n$,
$\operatorname{tr}(A)=\zeta+\zeta^{-1}=2\cos(2\pi/n)$ up to the choice of the
primitive root.  Thus $2\cos(2\pi/n)\in\Z$.  Since $-2\le 2\cos(2\pi/n)\le 2$,
the only possibilities are $2\cos(2\pi/n)\in\{-2,-1,0,1,2\}$, which occur
exactly at $n\in\{1,2,3,4,6\}$.
\end{proof}

\begin{corollary}
\label{cor:no32}
No planar lattice has an automorphism of order thirty-two; the largest
possible order of $\mathrm{Aut}(\Lam)$ is $12$, attained by the hexagonal
lattice (whose automorphism group is the dihedral group of order $12$).  In
particular, there is no ``order-thirty-two lattice torus'', so the second
clause of Conjecture~00000000277 is false.
\end{corollary}

\begin{proof}
By Theorem~\ref{thm:cryst} the rotational part of $\mathrm{Aut}(\Lam)$ is
contained in the rotations of order dividing $6$, and reflections are
involutions.  A finite subgroup of $O(2)$ generated by these is either cyclic
of order at most $6$ or dihedral of order at most $2\cdot 6=12$.  The bound
$12$ is achieved by the hexagonal lattice, whose automorphism group is the
dihedral group $D_6$ of order $12$.  In particular $32$ is not the order of
the automorphism group of any planar lattice.
\end{proof}

\section{The honest qualification: the first-eigenvalue reading}
\label{sec:honest}

We must distinguish two readings of $M$.

\begin{enumerate}
\item \emph{Literal reading (maximum over all eigenvalues).}  As in
Conjecture~00000000277, $M(T)$ is the maximum multiplicity of \emph{a}
Laplace eigenvalue, i.e.\ over all of $\operatorname{spec}(-\Delta_T)$.  Under
this reading Proposition~\ref{prop:unbounded} gives $M=\infty\neq6$: the claim
is false, and moreover no finite value can be correct.

\item \emph{First-eigenvalue reading (multiplicity of $\lambda_1$ only).}  If
instead $M$ were defined as the multiplicity of the \emph{first non-zero}
eigenvalue $\lambda_1$, then the value $6$ \emph{is} correct: it is attained by
the hexagonal (triangular) lattice, and the geometric argument (all shortest
dual vectors have equal length, pairwise angles at least $60^\circ$, and at
most six points on a circle can be $60^\circ$-separated) shows that no flat
torus has $\lambda_1$-multiplicity exceeding $6$.  This is the classical
statement that the two-dimensional kissing number is $6$.
\end{enumerate}

The conjecture is nevertheless false as written, for two separate reasons.

\begin{itemize}
\item Under the literal reading, $M=\infty$ (Proposition~\ref{prop:unbounded}),
so the asserted value $M=6$ is wrong.
\item Under \emph{both} readings, the clause ``multiplicity six is attained
only on the order-thirty-two lattice torus'' is impossible: by
Corollary~\ref{cor:no32} no planar lattice has automorphism group of order
$32$.  Even in the first-eigenvalue reading, where six is attainable, it is
attained by the hexagonal lattice with automorphism group of order $12$, not
by any order-thirty-two lattice.
\end{itemize}

Thus Conjecture~00000000277 is false as stated.

\begin{thebibliography}{9}
\bibitem{iwaniec}
H.~Iwaniec and E.~Kowalski,
\emph{Analytic Number Theory}, AMS Colloquium Publications \textbf{53},
American Mathematical Society, 2004.  (Jacobi's two-square theorem,
$\rtwo(N)=4(d_1(N)-d_3(N))$.)

\bibitem{conway-sloane}
J.~H.~Conway and N.~J.~A.~Sloane,
\emph{Sphere Packings, Lattices and Groups}, 3rd ed., Grundlehren der
mathematischen Wissenschaften \textbf{290}, Springer, 1999.
(Lattice automorphism groups, two-dimensional kissing number.)

\bibitem{berger}
M.~Berger,
\emph{G\'eom\'etrie}, vol.~1, CEDIC/Nathan, 1977.
(Crystallographic restriction: finite-order lattice automorphisms in the
plane have order in $\{1,2,3,4,6\}$.)

\bibitem{milnor}
J.~Milnor,
\emph{Eigenvalues of the Laplace operator on certain manifolds},
Proc.\ Nat.\ Acad.\ Sci.\ USA \textbf{51} (1964), 542.
(Spectra of flat tori and the dual-lattice description.)
\end{thebibliography}

\end{document}
